\documentclass[10pt,a4paper]{article}
\usepackage{amsmath,amssymb,amsthm,graphicx}
\usepackage{amsmath,amssymb,amsthm,graphicx}
\usepackage{hyperref}
\newtheorem{theorem}{Theorem}
\newtheorem{definition}[theorem]{Definition}
\newtheorem{proposition}[theorem]{Proposition}
\newtheorem{lemma}[theorem]{Lemma}
\newtheorem{claim}[theorem]{Claim}
\newtheorem{remark}[theorem]{Remark}
\newcommand{\C}{\mathbb C}
\newcommand{\N}{\mathbb N}
\newcommand{\fcom}{f_{\rm com}}
\newcommand{\Gcom}{\Gamma_{\rm com}}
\newcommand{\g}{\mathfrak g}
\newcommand{\End}{\mathrm{End}}
\newcommand{\GL}{\mathrm{GL}}
\newcommand{\ncw}{\mathrm{ncw}}
\DeclareMathOperator{\linspan}{linspan}
\DeclareMathOperator{\im}{im}
\DeclareMathOperator{\inset}{in}
\DeclareMathOperator{\outset}{out}
\allowdisplaybreaks
\AtBeginDocument{\author{Markus Bl\"aser, Christian Ikenmeyer, Meena Mahajan\\Anurag Pandey, Nitin Saurabh}}
\newcommand{\R}{\mathbb R}
\newcommand{\mw}{\mathrm{mw}}
\title{Monotone single-source algebraic branching programs are closed without changing their format}
\author{Markus Bläser, Christian Ikenmeyer, Meena Mahajan, Anurag Pandey, Nitin Saurabh}
\date{doi:10.4230/LIPIcs.CCC.2020.21; CC BY 3.0 Unported}
\begin{document}
\maketitle
\noindent\textit{Editorial note.} Editable transcription of complete formalCCBY3.0Unported Theorem3 proof and exactmodel definitions. No native retrieved or required. Original author affine/monotone/Laurentpositive conditions, normalization steps, termination andformat conclusion retained; no extra hypotheses or mathematical repair. Independent epsilon-elimination mechanism, not a trace nonclosure parameter variant. Actualfiveauthors metadata retained with two-line presentation. VQP/VNP observations andclass-level direct consequences not counted. Allproduction/validationYupeng; rootownscompilation.\par\medskip
\noindent\textit{Allowed standard prerequisites (editorial).}\begin{itemize}
\item ABP path semantics, affine linear labels and Laurent polynomials positive for smallpositiveepsilon. Exact formalSections1/4 model and full constructive authorproof retained; no prior nontrivial proof substituted for the epsilon-elimination process
\end{itemize}\medskip
\section*{Monotone single-(source,sink) ABPs and border width (original Sections1 and4)}
An ABP is a layered directed graph with $d+1$ layers of vertices (edges only go from layers $i$ to $i+1$). In the single-(source,sink) model the first and last layer have exactly one vertex. Each edge is labeled by a linear form, and the polynomial computed is $\sum_p\prod_{e\in p}\ell(e)$, where the sum is over paths from the source to the sink. The width is the number of vertices in its largest layer. In the affine setting labels are affine linear forms.
For the monotone setting let $R_d=\R_+[x_1,\ldots,x_m]_d$ and consider the affine or homogeneous setting. Since $\R$ is not algebraically closed, we switch to a more algebraic definition of approximation. Let $\R[\varepsilon,\varepsilon^{-1}]_+$ denote the ring of Laurent polynomials that are nonnegative for all sufficiently small $\varepsilon>0$. Clearly, elements from $\R[\varepsilon,\varepsilon^{-1}]_+$ can have a pole at $\varepsilon=0$ of arbitrarily high order. We define $\underline{\mw}(f)$ to be the smallest $w$ such that there exists a polynomial $f'$ over the ring $\R[\varepsilon,\varepsilon^{-1}]_+$ such that
\begin{itemize}
\item there exists a width-$w$ ABP over $\R[\varepsilon,\varepsilon^{-1}]_+$ that computes $f'$,
\item no coefficient in $f'$ contains an $\varepsilon$ with negative exponent, and setting $\varepsilon$ to0 in $f'$ yields $f$, i.e., $f'_{\varepsilon=0}=f$.
\end{itemize}
In the single-(source,sink) model the corresponding border and ordinary widths are denoted by $\underline{\mw}_1(f)$ and $\mw_1(f)$. The closure result gives $\underline{\mw}_1(f)=\mw_1(f)$ for all $f$.
Recall that an ABP has $d+1$ layers of vertices. If an ABP has $w_i$ many vertices in layer $i$, $1\le i\le d$, we say the ABP has format $w=(w_1,w_2,\ldots,w_d)$. We further recall that $w_{d+1}=w_1$.
\setcounter{theorem}{2}

\begin{theorem}\label{thm:monotone-closure}
Given a polynomial $f$ over $\R$ and given a format-$w$ single-(source,sink) ABP with affine linear labels over $\R[\varepsilon,\varepsilon^{-1}]_+$ computing $f_\varepsilon$ such that $\lim_{\varepsilon\to0}f_\varepsilon=f$. Then there exists a format-$w$ monotone single-(source,sink) ABP that computes $f$.
\end{theorem}

\section*{Complete author proof (original Section4, Theorem3)}
\begin{proof}
The proof is constructive and done by a two-step process. In the first step (which is fairly standard and works in many computational models) we move all the $\varepsilon$ with negative exponents to edges adjacent to the source. The second step then uses the monotonicity.
Given $\Gamma$ with affine linear labels over $\R[\varepsilon,\varepsilon^{-1}]_+$ we repeat the following process until all labels that contain an $\varepsilon$ with a negative exponent are incident to the source vertex.
\begin{itemize}
\item Let $e$ be an edge whose label contains $\varepsilon$ with a negative exponent $-i<0$. Moreover, assume that $e$ is not incident to the source vertex. Let $v$ be the start vertex of $e$. We rescale all edges outgoing of $v$ with $\varepsilon^i$ and we rescale all edges incoming to $v$ with $\varepsilon^{-i}$.
\end{itemize}
If we always choose the edge with the highest layer, then it is easy to see that this process terminates. Since every path from the source to the sink that goes through a vertex $v$ must use exactly one edge that goes into $v$ and exactly one edge that comes out of $v$, throughout the process the value of $\Gamma$ does not change. We finish this first phase by taking the highest negative power $i$ among all labels of edges that are incident to the source and then rescale all these edges with $\varepsilon^i$. The resulting ABP $\Gamma^i$ computes $\varepsilon^if_\epsilon$ and no label contains an $\varepsilon$ with negative exponent. We now start phase2 that transforms $\Gamma^i$ into $\Gamma^{i-1}$ that computes $\varepsilon^{i-1}f_\epsilon$ without introducing negative exponents of $\varepsilon$. We repeat phase2 until we reach $\Gamma^0$ in which we safely set $\varepsilon$ to0. Throughout the whole process we do not change the structure of the ABP and only rescale edge labels with powers of $\varepsilon$, which preserves monotonicity, so the proof is finished. It remains to show how $\Gamma^i$ can be transformed into $\Gamma^{i-1}$. An edge whose label is divisible by $\varepsilon$ is called an $\varepsilon$-edge. Consider the set $\Delta$ of vertices that are reachable from the source using only non $\varepsilon$-edges in $\Gamma^i$. The crucial insight is that since $\Gamma^i$ is monotone and computes a polynomial that is divisible by $\varepsilon$, we know that every path in $\Gamma^i$ from the source to the sink uses an $\varepsilon$-edge. Therefore $\Delta$ cannot contain the sink. We call a vertex in $\Delta$ whose outdegree is zero a leaf vertex. We repeat the following procedure until the source is the only leaf vertex:
\begin{itemize}
\item Let $v$ be a non-source leaf vertex in $\Delta$. We rescale all edges outgoing of $v$ with $\varepsilon^{-1}$ and we rescale all edges incoming to $v$ with $\varepsilon$.
\end{itemize}
It is easy to see that this process terminates with the source being the only leaf vertex. Since the source is a leaf vertex, all edges incident to the source are $\varepsilon$-edges. We divide all their labels by $\varepsilon$ to obtain $\Gamma^{i-1}$.
\end{proof}

\end{document}
